\section{Paths}\label{sec:11-path-algebras:02-paths:1}

A \textbf{path} in \(Q\) is a sequence of arrows that can be composed head-to-tail: \(\alpha_1, \alpha_2, \dots, \alpha_n\) with \(t(\alpha_i) = s(\alpha_{i+1})\) for each consecutive pair. In addition, for each vertex \(i\), a \textbf{trivial path} \(e_i\) of length \(0\) is allowed, starting and ending at \(i\) and composing with nothing but itself.

In the running example, the paths are: \(e_1, e_2, e_3\) (trivial), \(\alpha\), \(\beta\), and \(\beta\alpha\) (first \(\alpha\), then \(\beta\)), and nothing else, since there is no arrow out of \(3\) or into \(1\).

\subsection{References}\label{sec:11-path-algebras:02-paths:2}

Full citations in the \hyperref[sec:root:bibliography:1]{Bibliography}.

\begin{itemize}
\tightlist
\item
  Assem, Simson, and Skowroński (\cite{AssemSimsonSkowronski2006}), \textbf{Definition 1.2} (Chapter II, §1) --- path of length \(\geq 1\) as a sequence \((a \mid \alpha_1,\dots,\alpha_l \mid b)\), plus the trivial/stationary path \(\varepsilon_a\) of length \(0\) at each vertex \(a\).
\item
  Schiffler (\cite{Schiffler2014}), \textbf{Definition 2.1 and Example 2.2} (Chapter 2, §2.1) --- same notion, called the ``constant path'' (or ``lazy path'') \(e_i\) at vertex \(i\) for the length-\(0\) case.
\end{itemize}
